\documentclass[11pt]{amsart}
\usepackage{fullpage}

\usepackage[T1]{fontenc}
\usepackage[utf8]{inputenc}
\usepackage[slovene]{babel}
\usepackage{hyperref}
\usepackage{amsmath,amssymb,amsfonts,mathtools}

\linespread{1.1}
\setlength\parindent{6pt}

\newcommand{\N}{\mathbb{N}}
\newcommand{\Z}{\mathbb{Z}}
\newcommand{\R}{\mathbb{R}}
\newcommand{\E}{\mathbb{E}}
\newcommand{\Var}{\mathrm{Var}}
\newcommand{\Cov}{\mathrm{Cov}}
\newcommand{\Corr}{\mathrm{Corr}}
\newcommand{\WN}{\mathrm{WN}}
\newcommand{\IID}{\mathrm{IID}}
\newcommand{\AIC}{\mathrm{AIC}}
\newcommand{\BIC}{\mathrm{BIC}}

\theoremstyle{definition}
\newtheorem{definicija}{Definicija}[section]
\newtheorem{primer}[definicija]{Primer}
\newtheorem{opomba}[definicija]{Opomba}

\theoremstyle{plain}
\newtheorem{izrek}[definicija]{Izrek}
\newtheorem{trditev}[definicija]{Trditev}

\title{Analiza klimatskih spremenljivk z modeli \texttt{SARIMA} \\
\large{Pregled članka Dimri, Ahmad, Sharif (2020)}}
\author{[Janez Podlogar]}
\date{\today}

\begin{document}

\maketitle

\begin{abstract}

\end{abstract}

\section{Potrebno predznanje}

Ponovimo osnovne pojme in rezultate iz~\cite{basrak}, ki so potrebni za razumevanje članka.

\begin{definicija}
    \emph{Časovna vrsta} je stohastični proces $\{X_t\}_{t \in T}$, kjer je $T$ 
    indeksni množici, običajno $T = \Z$ ali $T = \N$.
\end{definicija}

\begin{definicija}
    Časovna vrsta $\{X_t\}_{t \in \Z}$ je \emph{šibko stacionarna}, če
    \begin{gather*}
        \forall t \in \Z. \exists m \in \R \colon \E[X_t] = m, \\ 
        \forall t \in \Z \colon \E[X_t^2] < \infty, \\
        \forall s,t,h \in \Z \colon \Cov(X_t, X_s) = \Cov(X_{t+h}, X_{s+h}).
    \end{gather*}
\end{definicija}

\begin{definicija}
    Naj bo $\{X_t\}_{t \in \Z}$ šibko stacionarna časovna vrsta. \emph{Avtokovariančna funkcija} je
    \begin{align*}
        \gamma_X \colon \Z &\to \R, \\
        h &\mapsto \Cov(X_{t}, X_{t+h}).
    \end{align*}
\end{definicija}

\begin{definicija}
    Naj bo $\{X_t\}_{t \in \Z}$ šibko stacionarna časovna vrsta. \emph{Avtokorelacijska funkcija}
    je
    \begin{align*}
        \rho_X \colon \Z &\to [-1,1], \\
        h &\mapsto \frac{\gamma_X(h)}{\gamma_X(0)}.
    \end{align*}
\end{definicija}


\begin{definicija}
    Šibko stacionarna časovna vrsta $\{X_t\}_{t \in \Z}$ je \emph{beli šum}, če
    \begin{gather*}
        \forall t \in \Z \colon \E[X_t] = 0, \\
        \gamma_X(h) = 
            \begin{cases}
                \sigma^2 &; h = 0, \\
                0 &; h \neq 0.
            \end{cases}
    \end{gather*}
    Beli šum je nekorelirano zaporedje z končno konstantno varianco in ničelno pričakovano vrednostjo.
    Označimo ga z $\WN(0, \sigma^2)$.
\end{definicija}

\begin{definicija}
    Naj bo $\{Z_t\}_{t \in \Z} \sim \WN(0,\sigma^2)$ in $\theta_1, \ldots, \theta_q \in \R$.
    Časovna vrsta $\{X_t\}_{t \in \Z}$, podana z
    \[
        X_t = Z_t + \theta_1 Z_{t-1} + \cdots + \theta_q Z_{t-q},
    \]
    je \emph{proces drsečega povprečja reda $q$}. Označimo ga z $\mathrm{MA}(q)$.
\end{definicija}

\begin{izrek}
    Proces $\mathrm{MA}(q)$ je šibko stacionaren.
\end{izrek}

\begin{proof}
    Ker je $\{Z_t\} \sim \WN(0, \sigma^2)$, za vsak $t \in \Z$ velja
    \[
        \E[X_t] = \E[Z_t + \theta_1 Z_{t-1} + \cdots + \theta_q Z_{t-q}] 
        = \sum_{j=0}^{q} \theta_j \E[Z_{t-j}] = 0,
    \]
    kjer označimo $\theta_0 = 1$.

    Za $h \geq 0$ in poljuben $t \in \Z$ velja
    \[
        \gamma_X(h) = \Cov(X_{t+h}, X_t) = \E[X_{t+h} X_t] 
        = \E\left[ \left(\sum_{i=0}^{q} \theta_i Z_{t+h-i}\right)
        \left(\sum_{j=0}^{q} \theta_j Z_{t-j}\right) \right].
    \]
    Zaradi linearnostjo pričakovane vrednosti dobimo
    \[
        \gamma_X(h) = \sum_{i=0}^{q} \sum_{j=0}^{q} \theta_i \theta_j \E[Z_{t+h-i} Z_{t-j}]
    \]
    in iz nekoleriranosti $\{Z_t\}$ sledi
    \[
        \E[Z_{t+h-i} Z_{t-j}] = 
            \begin{cases}
                \sigma^2 &; t+h-i = t-j,  \\
                0 &; \text{sicer}.
            \end{cases}
    \]
    Vsi čeni razen tist z $j = i - h$ so ničelni. Ker $0 \leq j \leq q$, velja
    mora veljati tudi $ h \leq i \leq g + h$. Za $0 \leq h \leq q$ dobimo
    \[
        \gamma_X(h) = \sum_{i=h}^{q} \theta_i \theta_{i-h} \sigma^2
                   = \sigma^2 \sum_{j=0}^{q-h} \theta_{j+h} \theta_j.
    \]
    Funkcijo $\gamma_X$ za $h < 0$ definiramo z $\gamma_X(h) = \gamma_X(-h)$, saj je
    avtokovariančna funkcija simetrična. Skupaj torej velja
    \[
        \gamma_X(h) = 
        \begin{cases}
            \sigma^2 \sum_{j=0}^{q-|h|} \theta_j \theta_{j+|h|} &; |h| \leq q, \\
            0 &; |h| > q.
        \end{cases}
    \]
    Ker je $\gamma_X(h)$ je odvisna le od $h$, ne pa od $t$, zato je proces šibko stacionaren.
\end{proof}

\begin{definicija}
    Naj bo $\{Z_t\}_{t \in \Z} \sim \WN(0,\sigma^2)$ in $\varphi_1, \ldots, \varphi_p \in \R$.
    Šibko stacionarna časovna vrsta $\{X_t\}_{t \in \Z}$, ki zadošča
    \[
        X_t = \varphi_1 X_{t-1} + \cdots + \varphi_p X_{t-p} + Z_t,
    \]
    je \emph{avtoregresijski proces reda $p$}. Označimo ga z $\mathrm{AR}(p)$.
\end{definicija}

\begin{thebibliography}{9}
    \bibitem{basrak} B.~Basrak, \emph{Essentials of Time Series with examples in R},
    Department of Mathematics, University of Zagreb.
\end{thebibliography}
d
\end{document}